\documentclass[10pt,a4paper]{amsart}
\usepackage[margin=19mm]{geometry}
\usepackage{amsmath,amssymb,mathtools}
\usepackage[scr=rsfs,cal=euler]{mathalfa}
\usepackage{xfrac}
\usepackage[unicode,hidelinks,bookmarks=false]{hyperref}
\newcommand{\ie}{i.e.\ }
\newcommand{\cf}{cf.\ }
\newcommand{\resp}{respectively\ }
\newcommand{\Subsection}{\subsection}
\providecommand{\nobreakdash}{\nobreak\hbox{-}}
\setlength{\emergencystretch}{2em}
\multlinegap0pt
\usepackage{bm}
\usepackage{bbm}
\usepackage{dsfont}
\usepackage{xspace}
\usepackage{cancel}
\usepackage{mathtools}
\usepackage{amsthm}
\makeatletter
\@ifundefined{dfn}{\newtheorem{dfn}{Dfn}}{}
\@ifundefined{lemd}{\newtheorem{lemd}[dfn]{Lemma}}{}
\makeatother
\newcommand{\Hom}{{\mathrm{Hom}}}
\newcommand{\RH}{{\mathrm {H}}}
\title[On the trilinear and Ginzburg-Rallis models]{On the trilinear and Ginzburg-Rallis models}
\author{Xinrui Wang}
\date{Source: arXiv:2606.31402v1 (CC BY 4.0)}
\begin{document}
\maketitle

\noindent\emph{Source and licence.} On the trilinear and Ginzburg-Rallis models. Version arXiv:2606.31402v1 (CC BY 4.0), distributed under the Creative Commons Attribution 4.0 International licence, as recorded in the arXiv licence metadata for this exact version and shown on the abstract page https://arxiv.org/abs/2606.31402v1. Full rights notice: licenses/arxiv-2606.31402-CC-BY-4.0.txt.\par\smallskip

\noindent\textit{Editorial scope.} Counted unit: the Lemd whose source label is \texttt{None} in the article (source0.tex lines 1342--1354, source label \texttt{None}), delivered together with the article's own complete original proof of it (source0.tex lines 1356--1451, one \texttt{proof} environment of 96 lines). No mathematics is written, repaired or re-derived: statement and proof are the article's own text line for line. Required context retained uncounted: none. Every definition, display and label of those lines is retained unchanged. Omitted from this package and not redistributed: 1--1341, 1355--1355, 1452--2076. Nothing inside a retained excerpt is omitted or altered: every retained body is delivered exactly as the article's lines are written, so no omission receipt of any kind accompanies this package. Citations: the retained excerpts carry no citation command at all, so no bibliography entry of the article is transcribed and no reference is redistributed. Reference adaptation: none was needed and none was made; every reference inside the delivered unit resolves inside it, so no unresolved reference or citation is produced. Printed slips kept literal: no printed word, symbol, hypothesis or number of the counted statement or of its proof was altered. Layout and class: the article's own class is \texttt{amsart} with packages including amsmath, geometry, amsthm, graphics, tabularx, amssymb, shapepar, eucal, verbatim, latexsym, amsfonts, amsxtra, bbm, color; the wrapper is the lane's pinned \texttt{amsart} editorial wrapper at 10pt on A4, which restates the theorem-like environments the retained text needs and adds the article's own macro definitions that the retained text needs (\texttt{\textbackslash Hom}, \texttt{\textbackslash RH}), with extra wrapper packages (bm, bbm, dsfont, xspace, cancel, mathtools). The delivered numbering is the unit's own. Assets: no retained span contains a figure environment, a graphics-inclusion command, a PDF-page inclusion, a table or a TikZ picture; no image file is copied into this package, the package declares no asset dependency and no image file is redistributed with it. Licence: version arXiv:2606.31402v1 of this article is distributed under the Creative Commons Attribution 4.0 International licence, as recorded in the arXiv licence metadata for this exact version and shown on the abstract page https://arxiv.org/abs/2606.31402v1. A full rights notice is delivered at licenses/arxiv-2606.31402-CC-BY-4.0.txt. Characters: every retained excerpt is pure ASCII, so the delivered engine, XeLaTeX with its default Latin Modern text font, reports no missing character.
\begin{lemd}
    Suppose that $ \gamma = (152643) $. If the space
    \[
        \Hom_{T_2}
        \left(
            J( \tau_3 )_\mathrm{ss}, ( \eta \circ \det ) \otimes \left( \omega_{ \tau_1 }^{-1} \otimes \omega_{ \tau_2 }^{-1} \right)
        \right)=0,
    \]
    then 
    \[
        \RH_i \left( P_\gamma, W_\gamma \right) = 0, \quad i \ge 0.
    \]
\end{lemd}

\begin{proof}
    Note that $T_\gamma'$ acts on $ J( \tau_1 ) \otimes J( \tau_2 ) \otimes J( \tau_3 ) $ by
    \[
        \tau_1 \left(
        \begin{bmatrix}
            g &  \\
              & g
        \end{bmatrix}
        \right) \otimes \tau_2 \left( 
        \begin{bmatrix}
            h &  \\
              & h
        \end{bmatrix}
        \right) \otimes \tau_3 \left(
        \begin{bmatrix}
            g &  \\
              & h
        \end{bmatrix}
        \right)
    \]
    and acts on $ \varphi_H^{ \gamma^{-1} } $ by $ \eta(gh) $. We have isomorphisms for all $ i \ge 0 $:
    \[
        \begin{aligned}
            \RH_i
            \left(
                T_\gamma', \RH_0 ( U_\gamma', W_\gamma )
            \right)
            \cong J_\psi ( \tau_1 ) 
            &\otimes J_\psi ( \tau_2 )\\
            &\otimes \RH_i
            \left(  
                T_2, J( \tau_3 ) \otimes \left( \omega_{ \tau_1 } \otimes \omega_{ \tau_2 } \right) \otimes \left( \eta^{-1} \circ \det \right) 
            \right).
        \end{aligned}
    \]
    Now consider the homology group of $T_2$. Arguing using the long exact sequence in homology, it is sufficient to show that the homology group remains zero if we replace $ J( \tau_3 ) $ by its semi-simplification $ J( \tau_3 )_\mathrm{ss} $. Write
    \[
        J( \tau_3 )_\mathrm{ss} = 
        \begin{cases}
            ( \xi_{11} \otimes \xi_{12} ) \oplus ( \xi_{21} \otimes \xi_{22} ), & \text{if } \tau_3 \text{ is a principal series representation,}\\
            \xi_1 \otimes \xi_2,                                                & \text{if } \tau_3 \text{ is a Steinberg representation,}       \\
            0,                                                                  & \text{if } \tau_3 \text{ is a supercuspidal representation.}
        \end{cases}
    \]
    In the case where $\tau_3$ is a principal series representation, we have
    \begin{equation}
        \label{ssspace}
        \begin{aligned}
             & \RH_i \left( T_2, J( \tau_3 )_\mathrm{ss} \otimes \left( \omega_{ \tau_1 } \otimes \omega_{ \tau_2 } \right) \otimes \left( \eta^{-1} \circ \det \right) \right)\\
            =& \bigoplus_{j = 1}^2 \RH_i \left( T_2, \left( \xi_{j1} \omega_{ \tau_1 } \eta^{-1} \right) \otimes \left( \xi_{j2} \omega_{ \tau_2 } \eta^{-1} \right) \right)   \\
            =& \bigoplus_{ \alpha + \beta = i } \bigoplus_{j = 1}^2 \RH_\alpha \left( k^\times, \xi_{j1} \omega_{ \tau_1 } \eta^{-1} \right) \otimes \RH_\beta \left( k^\times, \xi_{j2} \omega_{ \tau_2 } \eta^{-1} \right),
        \end{aligned}
    \end{equation}
    where the last equation follows from the K\"{u}nneth formula. Recall that for each character $\chi$ of $k^\times$,
    \[
        \RH_i ( k^\times, \chi ) =
        \begin{cases}
            \mathbb{C}, & \text{if } i = 0, 1 \text{ and } \chi = 1,\\
            0,          & \text{otherwise.}
        \end{cases}
    \]
    A straightforward calculation shows that
    \[
        ( \ref{ssspace} ) = 
        \begin{cases}
            \mathbb{C} \oplus \mathbb{C},                                     & \text{if } \xi_{l1} \omega_{ \tau_1 } = \xi_{l2} \omega_{ \tau_2 } = \eta \text{ and } i = 0, 2,                                                                                \\
            \mathbb{C} \oplus \mathbb{C} \oplus \mathbb{C} \oplus \mathbb{C}, & \text{if } \xi_{l1} \omega_{ \tau_1 } = \xi_{l2} \omega_{ \tau_2 } = \eta \text{ and } i = 1,\\
            0,                                                                & \text{otherwise.}
        \end{cases}
    \]
    Consequently, the vanishing of $ ( \ref{ssspace} ) $ for $ i = 0 $ guarantees its vanishing for all $ i \ge 0 $.
    Therefore, if $\tau_3$ is a principal series representation and 
    \[
        \Hom_{T_2}
        \left(
            J( \tau_3 )_\mathrm{ss},( \eta \circ \det ) \otimes \left( \omega_{ \tau_1 }^{-1} \otimes \omega_{ \tau_2 }^{-1} \right)
        \right)
        = 0,
    \]
    then by taking duality, we have
    \[
        \RH_0
        \left(
            T_2, J( \tau_3 )_\mathrm{ss} \otimes \left( \omega_{ \tau_1 } \otimes \omega_{ \tau_2 } \right) \otimes \left( \eta^{-1} \circ \det \right)
        \right)
        = 0.
    \]
    It follows that the spaces
    \[
        \RH_i
        \left(
            T_2, J( \tau_3 ) \otimes \left( \omega_{ \tau_1 } \otimes \omega_{ \tau_2 } \right) \otimes \left( \eta^{-1} \circ \det \right)
        \right)
    \]
    also vanish for all $ i \ge 0 $. By a similar argument, this assertion also holds when $\tau_3$ is a Steinberg representation, and holds automatically when $\tau_3$ is supercuspidal.
\end{proof}

\end{document}
